% ============================================================
% Constraint-First Topology: Foam Cobordisms, Nilpotent Transport,
% and Immersed Singularities as Admissibility Geometry
% ============================================================
\documentclass[12pt, a4paper]{article}

% Standard fonts for compatibility
\usepackage[T1]{fontenc}
\usepackage{mathpazo}
\usepackage{microtype}

% --- Math ---
\usepackage{amsmath, amsthm, amssymb, mathtools}
% unicode-math omitted for compatibility

% --- Layout ---
\usepackage[margin=1.15in]{geometry}
\usepackage{setspace}
\onehalfspacing
\usepackage{parskip}
\setlength{\parindent}{1.5em}
\setlength{\parskip}{0pt}

% --- Hyperlinks ---
\usepackage[
  colorlinks=true,
  linkcolor=black,
  citecolor=black,
  urlcolor=black
]{hyperref}

% Bibliography via thebibliography environment

% --- Theorem environments ---
\theoremstyle{plain}
\newtheorem{theorem}{Theorem}[section]
\newtheorem{proposition}[theorem]{Proposition}
\newtheorem{lemma}[theorem]{Lemma}
\newtheorem{corollary}[theorem]{Corollary}

\theoremstyle{definition}
\newtheorem{definition}[theorem]{Definition}
\newtheorem{example}[theorem]{Example}
\newtheorem{remark}[theorem]{Remark}

% Custom environments for interpretive claims
\newtheoremstyle{correspondence}
  {10pt}{10pt}{\itshape}{0pt}{\scshape}{.}{.5em}{}
\theoremstyle{correspondence}
\newtheorem{corrprin}[theorem]{Correspondence Principle}
\newtheorem{strucprin}[theorem]{Structural Equivalence Principle}
\newtheorem{interpprop}[theorem]{Interpretive Proposition}

% --- Operators and macros ---
\DeclareMathOperator{\Cob}{Cob}
\DeclareMathOperator{\AutIET}{Aut_{IET}}
\DeclareMathOperator{\SAF}{SAF}
\DeclareMathOperator{\Hom}{Hom}
\DeclareMathOperator{\End}{End}
\DeclareMathOperator{\im}{im}
\DeclareMathOperator{\rank}{rank}
\newcommand{\RR}{\mathbb{R}}
\newcommand{\QQ}{\mathbb{Q}}
\newcommand{\ZZ}{\mathbb{Z}}
\newcommand{\FF}{\mathbb{F}}
\newcommand{\BB}{\mathbb{B}}
\newcommand{\Kh}{\mathrm{Kh}}
\newcommand{\CLIO}{\mathrm{CLIO}}
\newcommand{\RSVP}{\mathrm{RSVP}}
\newcommand{\wedgeQ}{\wedge_{\QQ}}
\newcommand{\NN}{\mathcal{N}}
\newcommand{\calC}{\mathcal{C}}
\newcommand{\calX}{\mathcal{X}}
\newcommand{\calM}{\mathcal{M}}
\newcommand{\calA}{\mathcal{A}}
\newcommand{\calT}{\mathcal{T}}
\newcommand{\Tang}{\mathrm{Tang}}

% ============================================================
\begin{document}

\title{\large\scshape Constraint-First Topology:\\
\normalfont\large Foam Cobordisms, Nilpotent Transport,\\
and Immersed Singularities as Admissibility Geometry}

\author{\normalsize Flyxion \\ \small Independent Researcher}

\date{\small\today}

\maketitle

\begin{abstract}
Three recent papers in low-dimensional topology and homological algebra --- the
foam cobordism theory of Im and Khovanov, the nilpotent pair counting of Chen,
Im, Khovanov, Lillja, and Rugo, and the immersed surface functoriality of Carter,
Cooper, Khovanov, and Krushkal --- independently rediscover a recurring invariant
architecture that we term \emph{constraint-first process geometry}. In each domain,
meaningful structure does not reside in static objects but emerges as equivalence
classes of admissible local transformations under globally consistent constraint
systems. Weighted foam cobordisms encode conserved interval rearrangements whose
invariant lives in the exterior algebra $\RR \wedgeQ \RR$; nilpotent operator pairs
classify constrained bidirectional transport by depth-symmetric propagation modes
whose extinction behavior mirrors finite-time entropy collapse; and immersed surface
cobordisms with transverse double points extend categorical functoriality by
admitting controlled singularities as structured morphisms rather than pathological
failures of embedding.

We interpret all three constructions through the lens of the Relativistic Scalar-Vector
Plenum (RSVP) framework and the Constraint-Leveraged Inference Operator (CLIO),
arguing that these mathematical systems independently formalize three operational
faces of a single underlying architecture: process-primacy, obstruction-sensitivity,
and admissibility-governed invariant extraction. CLIO is introduced informally
through its operational role in each domain before being crystallized as a formal
operator in the synthesis. Our central claim is that constraint-first systems do not
derive coherence from objects; rather, objects emerge as stable equivalence classes
of admissible transformations.
\end{abstract}

\tableofcontents
\newpage

% ============================================================
\section{Introduction: The Recurring Architecture}
% ============================================================

There is a structural pattern appearing independently across several branches of
contemporary mathematics that is not yet widely recognized as a unified phenomenon.
The pattern concerns systems in which the fundamental entities are not objects but
transformations, and in which meaning is carried not by intrinsic properties of states
but by invariants preserved under classes of admissible local operations. Such systems
share three distinguishing features. First, they are \emph{process-primary}: the
ontologically basic entity is a transformation history or equivalence class thereof,
not an instantaneous configuration. Second, they are \emph{obstruction-sensitive}: the
invariants that survive constraint-preserving operations are precisely those that encode
what cannot be eliminated by any sequence of locally admissible moves. Third, they are
\emph{admissibility-governed}: coherence of the global structure is guaranteed not by
metric or embedding data but by the consistent closure of local rewrite or compatibility
laws.

Three papers in low-dimensional topology and linear algebra serve as our primary
mathematical sources. The foam cobordism paper of Im and Khovanov \cite{IK24foam}
develops a theory in which weighted oriented one-foams --- finite combinatorial graphs
whose edges carry positive real labels with conservation-compatible merge and split
rules at vertices --- are organized into cobordism equivalence classes. Their central
result identifies the cobordism group of such foams with the exterior product
$\RR \wedgeQ \RR$, the abelianization of the group of interval exchange transformations
via the Sah--Arnoux--Fathi (SAF) invariant, and the first K-theory group of the
Zakharevich assembler category. The identification is not merely algebraic: it reveals
that the combinatorial topology of weighted seams and cobordism moves is a faithful
encoding of interval exchange dynamics and its associated irreversibility obstruction.

The nilpotent pair paper of Chen, Im, Khovanov, Lillja, and Rugo \cite{CIKLR25} extends
a beautiful observation of Leinster concerning nilpotent operators on a single vector
space to pairs of linear maps between two spaces linked by bidirectional composition.
A pair $(f : V \to W,\, g : W \to V)$ is called nilpotent when the compositions $gf$
and $fg$ are both nilpotent. The paper introduces a key conceptual innovation --- the
distinction between \emph{balanced} and \emph{unbalanced} vectors --- which classifies
indecomposable nilpotent representations of the oriented two-cycle quiver. This
classification yields a closed counting formula over finite fields and a strikingly
elegant set-theoretic analogue connecting eventual constancy of map pairs with spanning
trees in complete bipartite graphs.

The immersed cobordism paper of Carter, Cooper, Khovanov, and Krushkal \cite{CCKK25}
extends Khovanov homology from smooth embedded surface cobordisms in four-space to
immersed surfaces possessing transverse double-point singularities. Rather than
treating double points as failures of embedding, the paper assigns specific chain maps
$A$ and $B$ to positive and negative nodes, extends the Carter--Saito movie move
calculus to include node interactions, and proves that the resulting 2-functor remains
well-defined on isotopy classes of immersed surfaces. The homological content of the
singular transitions is directly encoded by the Khovanov homology of the Hopf link.

Our purpose in this essay is not to compress these papers into a single formal system
but to exhibit a structural resonance among them that points toward a common underlying
architecture. We interpret this architecture through two complementary frameworks
developed in prior work: the \emph{Relativistic Scalar-Vector Plenum} (RSVP), a
field-theoretic model of admissibility-governed process evolution, and the
\emph{Constraint-Leveraged Inference Operator} (CLIO), an abstract operator class
governing constrained projection and trajectory compression. Both frameworks share the
central conviction that coherent structure --- physical, semantic, or topological ---
is not a property of configurations but of admissible transformations between them.

The essay proceeds in four parts. Part~I examines foam cobordisms as constraint-closure
invariants and the SAF obstruction as a topological irreversibility detector. Part~II
examines nilpotent pair transport as finite-time entropy exhaustion and balanced vectors
as coherence-preserving propagation modes. Part~III examines immersed surface cobordisms
as an extension of categorical morphism space admitting controlled singular transitions.
Part~IV synthesizes these correspondences and presents CLIO in its formal role, arguing
that all three mathematical systems independently instantiate the same constraint-first
process geometry.

Throughout, we distinguish three layers of structure that appear in each domain: local
transition laws governing individual moves, constraint-preserving global invariants
encoding obstruction information, and admissibility geometry determining which local
transformations preserve coherent global structure. This tripartite organization is not
imposed but discovered, appearing in each of the three source papers nearly
uncannily well.

% ============================================================
\section{Part I: Foam Cobordisms as Constraint-Closure Invariants}
% ============================================================

\subsection{Weighted foams and conserved interval decomposition}

A weighted oriented one-foam in the sense of Im and Khovanov is a finite oriented
trivalent graph whose edges carry positive real labels, subject to the conservation
condition that at every merge or split vertex the labels sum consistently: if two edges
of weights $a$ and $b$ merge into a single edge, that edge carries weight $a + b$.
This is not a metric constraint. The number $a$ does not measure the geometric length
of an edge in any embedding space. It is a \emph{compatibility label} ensuring that
the vertex interaction respects a local conservation law. The paper is explicit about
this: the thickening constructions in Sections~2.4.1 and 2.4.2 of \cite{IK24foam}
interpret these labels as measures of interval width in the lower-limit topology or as
weights in a measured foliation, but the algebraic structure is entirely independent
of any embedding.

This is the first appearance of what we will call \emph{constraint-before-content}: the
label $a$ is meaningful not as an intrinsic property of an edge in isolation but only
relative to the compatibility constraints imposed at vertices. An edge label acquires
significance through its participation in admissible merge and split events.

Formally, given a commutative semigroup $(H, +)$, one defines the cobordism group
$\Cob^1_H$ of $H$-weighted oriented one-foams modulo $H$-weighted oriented two-foam
cobordisms. The main theorem of Section~2 of \cite{IK24foam} then asserts:

\begin{theorem}[Im--Khovanov \cite{IK24foam}]
\label{thm:IK-main}
There are isomorphisms of abelian groups
\[
\Cob^1_{>0} \;\cong\; \RR \wedgeQ \RR \;\cong\; H_1(\AutIET, \ZZ) \;\cong\; K_1(\calC_Z),
\]
where $\Cob^1_{>0}$ denotes the cobordism group of weighted oriented one-foams with
positive real labels, $\AutIET$ is the group of interval exchange transformations of
$[0,1)$, $\calC_Z$ is the Zakharevich assembler category for interval exchange
transformations, and the second isomorphism is the Sah--Arnoux--Fathi invariant.
\end{theorem}

The proof proceeds by constructing an isomorphism $\psi : \Cob^1_{>0} \xrightarrow{\sim} \RR \wedgeQ \RR$
via a local invariant $\nu$ assigned to planar projections of one-foams. The invariant
accumulates contributions from merge and split vertices according to their thin-edge
orderings and from virtual crossings between interval strands. One key calculation is
that the invariance of $\nu$ under move~1 in Figure~15 of \cite{IK24foam} requires
exactly the relation $a \wedge b + b \wedge a = 0$, and invariance under move~5 requires
the bilinearity of the tensor product. The cobordism relations are thus the geometric
shadows of the algebraic relations defining $\RR \wedgeQ \RR$.

\subsection{The SAF invariant as topological irreversibility}

The Sah--Arnoux--Fathi invariant of an interval exchange transformation
$T_{\lambda, \sigma}$ is defined as
\[
\SAF(T_{\lambda,\sigma}) = \sum_{i=1}^r \lambda_i \otimes t_i
= 2 \sum_{i < j \,:\, \sigma(j) < \sigma(i)} \lambda_i \wedge \lambda_j,
\]
where $t_i$ is the displacement of the $i$-th interval under $\sigma$. It vanishes
precisely on the commutator subgroup of $\AutIET$, yielding the short exact sequence
\[
1 \to [\AutIET, \AutIET] \to \AutIET \xrightarrow{\SAF} \RR \wedgeQ \RR \to 1.
\]
The invariant detects the \emph{irremovable asymmetry} of an interval permutation:
those rearrangements that cannot be expressed as products of commutators are precisely
those whose SAF class is nontrivial in the exterior algebra.

From the perspective of RSVP admissibility geometry, this is structurally very close to
a topological entropy circulation invariant. The system locally permutes intervals without
altering total measure, yet globally there persists a conserved antisymmetric quantity
that survives all locally admissible rearrangements. In RSVP, similarly, field
configurations related by admissible entropy-compatible evolution share the same
effective obstruction class, and what distinguishes genuine physical trajectories from
arbitrary rearrangements of configuration data is the preservation of such circulation
invariants.

\begin{corrprin}[SAF Obstruction and Constraint-Closure Geometry]
\label{cp:SAF}
The Sah--Arnoux--Fathi invariant of an interval exchange transformation is the minimal
admissibility invariant of interval rearrangement dynamics: it encodes exactly that
asymmetric relational structure which no sequence of locally admissible reassociations
can eliminate. In RSVP-style admissibility geometry, this corresponds to the role of
topological obstruction classes in trajectory compression: the SAF class is the
residual entropy-circulation datum surviving constraint-preserving coarse-graining.
\end{corrprin}

The wedge structure $\lambda_i \wedge \lambda_j$ is especially significant. It is
antisymmetric in the pair $(\lambda_i, \lambda_j)$, so it records the \emph{order} of
the two intervals relative to the permutation, not merely their coexistence. This is not
metric information. It is \emph{relational orientation information} surviving the
abelianization: the class $\lambda_i \wedge \lambda_j \in \RR \wedgeQ \RR$ witnesses
the fact that $i$ and $j$ were inverted by $\sigma$ in a way that no local cobordism
can undo.

\subsection{The antisymmetric bracket and local reassociation laws}

Section~3 of \cite{IK24foam} develops a parallel theory for planar unoriented weighted
one-foams. The cobordism group $\Cob^{1,\mathrm{up}}_{>0}$ of such foams is generated
by tripod foams $T(a,b)$ and is isomorphic to the group $Z^2(\RR_{>0})$ defined by
generators $[a,b]$ with $a,b > 0$ and the relations
\begin{align}
[a,a] &= 0, \label{eq:antisym1} \\
[a,b] + [b,a] &= 0, \label{eq:antisym2} \\
[a,b] + [a+b,c] &= [a,b+c] + [b,c]. \label{eq:cocycle}
\end{align}
Relation~\eqref{eq:cocycle} is the defining identity of this algebraic structure. Im and
Khovanov observe that it is a two-cocycle relation: it can be interpreted as the signed
boundary of a three-simplex with oriented edges labeled by $a$, $b$, $c$, $a+b$, $b+c$,
$a+b+c$. Geometrically, it expresses the commutativity of two distinct sequences of
vertex cobordisms merging three parallel strands of weights $a$, $b$, $c$ into a single
strand of weight $a+b+c$.

From our perspective, relation~\eqref{eq:cocycle} is a \emph{local reassociation law}.
It asserts that two distinct orderings of a three-way merge are equivalent in the
cobordism group. This is formally very close to the Mac Lane pentagon coherence
condition in monoidal category theory, and it resembles the entropy redistribution
identities appearing in TARTAN recursive tiling: the decomposition order changes, but
the global constraint-compatible invariant does not.

\begin{strucprin}[Reassociation Invariance and TARTAN Constraint Closure]
\label{sp:tartan}
The antisymmetric two-cocycle relation \eqref{eq:cocycle} is the foam-theoretic
instance of a general reassociation invariance principle: any two admissible
sequences of local decomposition events producing the same boundary configuration
are equivalent under constraint-consistent cobordism. In the TARTAN framework, an
analogous principle governs recursive tiling and constraint propagation: distinct
local rewrite sequences preserving boundary admissibility conditions are identified
in the global constraint-closure quotient.
\end{strucprin}

The bracket $[a,b]$ is \emph{almost} bilinear: Proposition~3.5 of \cite{IK24foam} shows
that $2([a, b_1 + b_2] - [a,b_1] - [a,b_2]) = 0$ in $Z^2(\RR_{>0})$. The deviation from
linearity is a two-torsion obstruction. In RSVP terms, this two-torsion corresponds
to a residual orientational ambiguity in the constraint-closure quotient: the system
is almost, but not quite, indifferent to linear decomposition order.

\subsection{Cobordism classes as admissibility equivalences}

The most conceptually important move in \cite{IK24foam} is the replacement of geometric
identity with cobordism equivalence. Two weighted one-foams are declared equivalent not
when they are homeomorphic as graphs with labels but when one can be transformed into
the other through a sequence of elementary two-foam cobordisms: cap and cup moves
adding or removing circles, saddle cobordisms, vertex cobordisms passing through a
seam interaction, and isotopies of planar diagrams. Each elementary cobordism preserves
the local conservation law at seams. The global equivalence class encodes exactly what
cannot be undone by any sequence of such moves.

This is the foam-theoretic realization of what we mean by a \emph{constraint-first
ontology}. The meaningful object is not the weighted graph itself but its equivalence
class under admissible local operations. The operations are primary; the equivalence
class --- the cobordism group element --- is derived. Obstruction to triviality is
encoded by the SAF invariant or the bracket class $[a,b]$, and these are precisely
the relational data that survive all locally admissible transformations.

In the lower-limit-topology thickening of Section~2.4.1 of \cite{IK24foam}, a
weighted one-foam $U$ is replaced by a topological space $T(U)$ locally homeomorphic
to $[0,1)_\ell \times (0,1)$, where $[0,1)_\ell$ carries the lower-limit topology. The
foliation of $T(U)$ by leaves $x \times (0,1)$ gives a measured foliation in which the
leaf-compactness condition (all leaves being closed circles) is exactly the condition
for $U$ to be null-cobordant. This is a local-to-global theorem mediated by
admissibility: what is geometrically local (leaf compactness) determines what is
globally trivial (cobordism class).

% ============================================================
\section{Part II: Nilpotent Transport and Finite-Time Entropy Exhaustion}
% ============================================================

\subsection{Bidirectional transport and quiver representations}

The paper \cite{CIKLR25} studies pairs of linear maps $(f : V \to W,\, g : W \to V)$
between finite-dimensional vector spaces, asking when the composed endomorphisms
$T = gf : V \to V$ and $T' = fg : W \to W$ are both nilpotent. This is a representation-
theoretic question: such a pair is precisely a representation of the oriented two-cycle
quiver
\[
\Gamma \;=\; \bigl( V \underset{g}{\overset{f}{\rightleftharpoons}} W \bigr),
\]
and nilpotency of the compositions is the nilpotency condition on the representation.

The shift from a single endomorphism to a pair of maps in opposite directions is
mathematically significant. For a single nilpotent endomorphism $T : V \to V$, the
indecomposable representations are Jordan blocks $J_r$ of varying sizes $r \geq 1$,
each characterized by a single integer. For the two-cycle quiver, indecomposable
nilpotent representations are classified by a pair of integers $(m,n)$ with $|m-n| \leq 1$,
reflecting the bidirectional transport structure.

From the RSVP perspective, this shift models the transition from a single state manifold
with self-referential evolution to a coupled pair of manifolds linked by transport
operators. The compositions $T = gf$ and $T' = fg$ are not the transports themselves
but the \emph{round-trip operators} measuring what survives a complete transit in both
directions. Nilpotency of $T$ means that iterated round-trips eventually annihilate all
state vectors: the reachable future volume under repeated bidirectional transport
shrinks to zero in finite time.

\subsection{Balanced and unbalanced vectors as propagation modes}

The central conceptual innovation of \cite{CIKLR25} is the distinction between balanced
and unbalanced vectors. Given a nilpotent pair $(f,g)$ and a vector $v \in V$, one
forms two ascending chains of iterates:
\[
v, Tv, T^2 v, \ldots, T^{a-1}v \neq 0, \quad T^a v = 0
\]
\[
fv, T'fv, (T')^2 fv, \ldots, (T')^{\ell-1}fv \neq 0, \quad (T')^\ell fv = 0.
\]
The vector $v$ is \emph{balanced} if $a = \ell$, meaning the two chains have equal
length and $f$ restricts to an isomorphism $T[v] \xrightarrow{\sim} T'[fv]$ between the
corresponding cyclic subspaces. It is \emph{unbalanced} if $\ell = a - 1$, meaning the
image chain is shorter by exactly one step.

This dichotomy classifies the propagation modes of the bidirectional transport system.
A balanced vector is a mode in which information propagates symmetrically through the
coupled system: the depth of recursive iteration is equal in both spaces, and the
transport map $f$ identifies the two recursive histories. An unbalanced vector is a
dissipative mode: information propagates one step further in the source space $V$ than
in the target space $W$, producing an asymmetric accessibility profile.

\begin{corrprin}[Balanced Propagation and Coherent Transport]
\label{cp:balanced}
Balanced vectors in a nilpotent pair $(f,g)$ are the constraint-preserving propagation
modes of bidirectional transport: they witness symmetric recursive iteration depth
across the coupled manifolds $V$ and $W$ and correspond to coherent transport channels
in which no dimensional collapse asymmetry arises. Unbalanced vectors are dissipative
modes encoding one-sided information loss. In RSVP admissibility geometry, this
distinction corresponds to the decomposition of transport operators into
coherence-preserving and entropy-producing sectors.
\end{corrprin}

The paper's generalization of Leinster's bijection is the following result.

\begin{theorem}[Chen--Im--Khovanov--Lillja--Rugo \cite{CIKLR25}]
\label{thm:CIKLR-main}
There is a bijection
\[
\Hom(V,W) \times \Hom(W,V) \times (V \cup_0 W) \;\cong\; \NN(V,W) \times V \times W,
\]
where $\NN(V,W)$ denotes the set of nilpotent pairs and $V \cup_0 W$ is the union of
$V$ and $W$ along their zero vectors. Over a finite field $\mathbb{F}_q$, this yields
the closed formula
\[
N_{m,n} = q^{2mn - m - n}(q^m + q^n - 1)
\]
for the number of nilpotent pairs with $\dim V = m$ and $\dim W = n$.
\end{theorem}

The proof passes through Lemma~3.14 of \cite{CIKLR25}, which constructs an explicit
bijection between quadruples with balanced components and quadruples with unbalanced
components by a controlled redefinition of $f$ on the subspace $T[v_b]$. This bijection
is the algebraic heart of the paper and is structurally reminiscent of a gauge
transformation: the pair $(f,g)$ is modified on a single subspace to exchange the
balanced and unbalanced sectors without changing the global nilpotency class.

\subsection{Nilpotency as finite-time entropy exhaustion}

The nilpotent condition $T^k = 0$ for some finite $k$ has a clear geometric
interpretation in RSVP language. Define the \emph{accessibility volume} of the
transport system at depth $n$ as the dimension of $\im(T^n)$. Nilpotency asserts that
this volume decreases strictly and reaches zero at the finite depth $k$. The sequence
\[
\dim V = \dim \im(T^0) \geq \dim \im(T^1) \geq \cdots \geq \dim \im(T^k) = 0
\]
is an exact analogue of entropy descent: the logarithm of the accessibility volume
decreases monotonically to $-\infty$ under iterated application of $T$.

In RSVP, entropy is defined geometrically as the logarithmic accessible future volume
in trajectory space under admissibility constraints. A nilpotent transport operator is
therefore a \emph{finite-time accessibility extinguisher}: a morphism of the state
manifold under which all admissible future structure is annihilated within $k$ steps.

\begin{strucprin}[Nilpotent Exhaustion and RSVP Entropy Descent]
\label{sp:nilp}
A nilpotent pair $(f,g)$ with nilpotency degree $k$ is the algebraic model of a
finite-time entropy exhaustion process: the accessibility volume of the bidirectional
transport system descends strictly to zero in exactly $k$ iterations of the round-trip
operator. This is structurally equivalent to a bounded-depth admissibility collapse in
RSVP trajectory geometry, in which the admissible future cone shrinks to a point in
finite proper time under constraint-preserving evolution.
\end{strucprin}

\subsection{The set-theoretic model and acyclic causal structure}

The set-theoretic analogue developed in Section~2 of \cite{CIKLR25} is remarkably
clean. For finite sets $X$ and $Y$ with $|X| = m$ and $|Y| = n$, a pair of maps
$(f : X \to Y, g : Y \to X)$ is \emph{eventually constant} when the compositions $(gf)^k$
and $(fg)^k$ eventually map all elements to a unique 2-cycle attractor $(x_0, y_0)$
with $f(x_0) = y_0$ and $g(y_0) = x_0$.

Theorem~2.4 of \cite{CIKLR25} establishes a bijection between eventually constant
pairs and spanning trees of the complete bipartite graph $K(m,n)$ together with a
choice of distinguished edge. The counting formula is $m^{n-1} n^{m-1}(m+n-1)$.

The spanning-tree correspondence is conceptually central. A spanning tree on a
bipartite graph is a minimal connected acyclic structure: it contains all vertices,
preserves global reachability, and eliminates all cycles. In RSVP and CLIO terms,
trees are \emph{minimal admissibility skeletons}: they retain just enough connectivity
to support constrained propagation while removing all redundant entropy-producing
feedback.

The eventually constant condition is precisely acyclicity of the causal structure
outside the minimal 2-cycle attractor. All state transitions flow toward the attractor
along tree-like dependency paths. This mirrors the TARTAN principle that stable
recursive systems enforce directed acyclic propagation and eliminate self-referential
closure loops.

\subsection{Boolean nilpotency and directed acyclic causal graphs}

The Boolean semiring $\BB = \{0,1 : 1 + 1 = 1\}$ provides the sharpest version of the
acyclicity correspondence. Over $\BB$, Proposition~4.2 of \cite{CIKLR25} shows that a
Boolean matrix $A \in \mathrm{Mat}_n(\BB)$ is nilpotent if and only if the directed
graph with adjacency matrix $A$ is a directed acyclic graph (DAG). The proof is
immediate: a $1$ on the diagonal of any power $A^k$ is the algebraic signature of an
oriented cycle, and nilpotency is exactly the condition excluding all such cycles.

\begin{interpprop}[Boolean Nilpotency and Causal Admissibility]
\label{ip:boolean}
Nilpotent Boolean matrices are in bijection with directed acyclic graphs on $n$ ordered
vertices. The nilpotency condition is therefore equivalent to the absence of any
self-reinforcing causal cycle in the associated computation graph. In the context of
CLIO and RSVP, this formalizes the principle that constraint-admissible recursive
computation is acyclic: admissible inference propagation proceeds without
pathological self-reference, and entropy-respecting modular composition excludes
runaway feedback loops.
\end{interpprop}

The recurrence $a_n = \sum_{k=1}^n (-1)^{k-1} \binom{n}{k} 2^{k(n-k)} a_{n-k}$
counting DAGs on $n$ ordered vertices (Corollary~4.3 of \cite{CIKLR25}) encodes the
combinatorial complexity of admissible causal structures. Each term in the inclusion-
exclusion alternates signs in the manner of an obstruction-theoretic count: positive
contributions from acyclic structures, negative corrections removing configurations
containing at least one cycle.

% ============================================================
\section{Part III: Immersed Cobordisms as Controlled Singular Transitions}
% ============================================================

\subsection{Extending the morphism category}

The paper \cite{CCKK25} addresses a fundamental question in categorical topology:
how far can functoriality be extended? Classical Khovanov homology associates to each
oriented link $L \subset \RR^3$ a bigraded abelian group $\Kh(L)$ and to each smooth
oriented surface cobordism $\Sigma \subset \RR^3 \times [0,1]$ between links a graded
map $\Kh(L_0) \to \Kh(L_1)$, well-defined up to overall sign on isotopy classes of
surfaces. The 2-functor
\[
\kappa : \Tang^4_b \to K_p
\]
sends balanced tangle cobordisms to chain homotopy classes of maps between Khovanov
chain complexes.

The extension problem is: what happens when the surface $\Sigma$ is no longer smoothly
embedded but merely immersed, possessing transverse self-intersection points modeled
locally by the complex variety $\{(w,z) \in \mathbb{C}^2 : wz = 0\}$? This local model
is the intersection of two complex planes in $\mathbb{C}^2$, whose link in a small
3-sphere is a Hopf link.

Rather than declaring such immersions inadmissible, Carter, Cooper, Khovanov, and
Krushkal define chain maps $A$ and $B$ associated to positive and negative double-point
singularities respectively and prove that the resulting assignment extends $\kappa$
to a 2-functor
\[
\tilde\kappa : \Tang^4_{\hat{b}} \to K_p
\]
on the larger category $\Tang^4_{\hat{b}}$ of balanced immersed cobordisms with double
points.

This is a categorical expansion of admissible morphisms: the paper does not embed
singular surfaces into a pre-existing smooth theory but genuinely enlarges the 2-category
of cobordisms by adding morphism types that carry new algebraic information.

\subsection{Singularities as structured computational events}

The maps $A : C^+ \to C^-$ and $B : C^- \to C^+$ assigned to double points are
explicitly computed from the chain complexes associated to positive and negative
crossings. The map $A$ has homological degree $-2$ (in the $t$-grading) and lives in
the class of $(t,q)$-degree $(-2,-6)$ corresponding to a generator of the
$q$-shifted negative Hopf link homology $q^{-2} H^-$. The map $B$ has degree $+2$
and corresponds to a generator of $q^{-2} H^+$.

The appearance of Hopf link homology here is structurally significant. The Hopf link
is the link of the double-point singularity: it bounds a small 3-sphere around the
node in 4-space, and its Khovanov homology records the algebraic topology of that
local singular region. The maps $A$ and $B$ are therefore not arbitrary choices but
are determined by the intrinsic algebraic topology of the singularity type.

\begin{interpprop}[Double Points as Singular Morphism Types]
\label{ip:double}
Transverse double-point singularities in immersed surface cobordisms are not
pathological failures of the theory but structured morphism types carrying definite
homological information: the chain maps $A$ and $B$ encode the Khovanov homology of
the Hopf link, which is the canonical link of the double-point singularity. In RSVP
admissibility geometry, this corresponds to the role of entropy concentration loci as
structured transition operators rather than excluded defects: a singularity is an
admissible event carrying definite obstruction data.
\end{interpprop}

The homological degree shift is important. Ordinary embedded cobordisms produce
degree-zero maps between Khovanov homology groups (after appropriate grading shifts).
The singular maps $A$ and $B$ shift the $t$-degree by $\pm 2$, allowing access to
homological classes that are not reachable by smooth embedded cobordisms. The paper
explicitly frames this as a probe of ``non-geometric'' homology: the classes in
$\Kh^{i,j}(L)$ that cannot be represented as images of smooth surface maps but can
be reached through controlled singular transitions.

In RSVP language, this corresponds to the distinction between low-entropy embedded
trajectories (which access only a narrow sector of the accessible state space) and
controlled entropy-transition operators (which allow passage between sectors separated
by admissibility barriers).

\subsection{Movie moves as rewrite grammar over admissible transitions}

The Carter--Saito movie move framework, extended in Section~4 of \cite{CCKK25}, is a
formal grammar for surface isotopies in four-dimensional space. A movie of a surface
$F \subset \RR^3 \times [0,1]$ is a sequence of cross-sectional stills --- oriented
links in $\RR^3 \times \{s_i\}$ --- between which a finite list of elementary events
occurs: births and deaths of circles, saddle moves, Reidemeister moves of types I, II,
and III, and, in the extended theory, double-point creation or annihilation events
(nodes).

The movie move theorem (Theorem~4.1 of \cite{CCKK25}) asserts that any two movies
presenting isotopic immersed surfaces are related by a finite sequence of movie moves
from the extended list MM1--MM19. This is a \emph{completeness theorem for the rewrite
grammar}: any two equivalent representations of the same isotopy class are connected
by finitely many elementary substitutions.

From the perspective of TARTAN and RSVP, this is the topological realization of what
we call \emph{trajectory-space equivalence}: the surface is not a static object but an
equivalence class of local transition histories, and the movie moves are the rewrite
rules identifying different local presentations of the same higher-dimensional process.
Different movies are different \emph{semantic foliations} of the same underlying
admissibility structure.

\begin{strucprin}[Movie Moves as Gauge Invariance over Presentation]
\label{sp:movies}
The extended Carter--Saito movie moves function as gauge equivalences over local
presentations of immersed surface cobordisms: any two movies presenting the same
isotopy class of immersed surface are related by finitely many elementary rewrite
substitutions. In the RSVP and TARTAN frameworks, this is the categorical instance
of a general invariance principle: the meaningful entity is not the local presentation
(the movie) but the equivalence class of presentations under admissible rewrite
operations, and coherence of the global structure is guaranteed by the closure of
the rewrite algebra.
\end{strucprin}

\subsection{The chart calculus and projection geometry}

The chart formalism developed in Section~4.2 of \cite{CCKK25} provides a two-dimensional
encoded representation of the full four-dimensional immersed surface. A chart is an
oriented labeled graph in the $(s,y)$-plane encoding, via decorated edges and vertices,
the fold singularities, crossing arcs, cusps, triple points, and node events of the
surface projected into three-dimensional space.

Formally, the chart is a structured compression of the four-dimensional process into a
two-dimensional combinatorial object retaining all information necessary to reconstruct
admissible transformations. This is a direct analogue of the CLIO projection
\[
\pi : \calX \to \calM
\]
where $\calX$ is the full trajectory space of the surface in $\RR^3 \times [0,1]$ and
$\calM$ is the compressed representation space carrying the invariant algebraic data.
The chart does not record the surface's geometry; it records the admissibility structure
of its transformations.

The bookmark code --- the combinatorial encoding of a cross-sectional still via symbols
for crossings, maxima, minima, and vertical strands --- functions as a local coordinate
system in $\calM$. Different bookmark words for the same link diagram are related by
Reidemeister moves, just as different local coordinates on $\calM$ are related by the
gauge equivalences of the theory.

\subsection{Reaching non-geometric homology through singular transitions}

A recurrent theme in the topology of four-manifolds is that smooth embedding severely
restricts the algebraic invariants accessible to geometric probes. The Whitney trick
fails in dimension four: two embedded surfaces in a 4-manifold with algebraically
canceling intersection points cannot in general be isotoped apart, and the obstruction
is not purely algebraic but depends on the smooth structure. Double points that cannot
be removed by homotopy are a fundamental feature of four-dimensional topology rather
than an exceptional pathology.

Carter, Cooper, Khovanov, and Krushkal observe that ordinary embedded cobordisms
produce maps concentrated near the diagonal of Khovanov homology (along the line
$q = 2t + c$ for fixed constants $c$ depending on grading conventions), whereas the
singular maps $A$ and $B$ have degrees $(-2,-6)$ and $(2,4)$ respectively, accessing
off-diagonal regions not reachable by smooth cobordisms. Remark~3.4 of \cite{CCKK25}
frames this as the central motivation: the extension $\Tang^4_b \subset \Tang^4_{\hat{b}}$
can be understood as the simplest nontrivial enlargement of the theory along homology
classes not represented geometrically.

In RSVP terms: smooth embedded trajectories are \emph{low-entropy morphisms} accessing
only a narrow sector of the reachable state space. Controlled singular transitions are
\emph{entropy-transition morphisms} crossing admissibility barriers that smooth
embeddings cannot traverse. The full accessibility map of Khovanov homology is larger
than what smooth geometry alone can explore.

% ============================================================
\section{Part IV: Synthesis and the CLIO Operator}
% ============================================================

\subsection{The tripartite architecture across three domains}

We are now in a position to identify the common architecture appearing in all three
source papers. In each domain, there is a precise tripartite organization of the
mathematical structure into local transition laws, global constraint invariants, and
admissibility geometry. We exhibit this organization explicitly.

In foam cobordism theory, the \emph{local transition laws} are the elementary cobordisms
of Figure~16 in \cite{IK24foam}: cap and cup moves, saddle cobordisms, vertex
cobordisms crossing a seam interaction, and isotopies of planar foam diagrams. These
are the primitive admissible local events. The \emph{global constraint invariants} are
the SAF class $\SAF(T) \in \RR \wedgeQ \RR$ and the antisymmetric bracket class
$[a,b] \in Z^2(\RR_{>0})$, both preserved under all sequences of elementary cobordisms.
The \emph{admissibility geometry} is the weighted conservation law at foam vertices:
$a + b \leftrightarrow (a,b)$ is the local compatibility constraint determining which
elementary cobordisms are permitted.

In nilpotent transport theory, the \emph{local transition laws} are the block-matrix
redefinitions of $f$ on cyclic subspaces (as in Lemma~3.14 of \cite{CIKLR25}), the
Fitting decomposition into stable and nilpotent sectors (Lemma~3.3), and the bijective
rebalancing between balanced and unbalanced vector sectors. The \emph{global constraint
invariants} are the nilpotency degree $k$ of the round-trip operators $T = gf$ and
$T' = fg$, the dimension pair $(\dim V_I, \dim V_N)$ of the Fitting decomposition, and
the counting invariant $N_{m,n} = q^{2mn-m-n}(q^m + q^n - 1)$. The \emph{admissibility
geometry} is the nilpotency condition itself: $T^k = 0$ is the constraint filtering
which pairs $(f,g)$ are admissible nilpotent transport systems.

In immersed cobordism theory, the \emph{local transition laws} are the elementary movie
events: births, deaths, saddles, Reidemeister moves of types I--III, $\psi$-moves, and
the new double-point node events encoded by $A$ and $B$. The \emph{global constraint
invariants} are the Khovanov homology groups $\Kh(L)$ of the boundary links and the
homotopy classes of chain maps between them. The \emph{admissibility geometry} is the
coherence conditions ensuring that all extended movie moves commute up to chain
homotopy: this is the constraint guaranteeing that the assignment of chain maps to
movies is well-defined on isotopy classes.

\subsection{Process primacy as a mathematical principle}

The deepest shared feature of the three theories is that in each case the fundamental
object is a \emph{transformation class} rather than a \emph{configuration}. A weighted
one-foam is not primarily a graph with labels; it is a class of configurations related
by admissible elementary cobordisms. A nilpotent pair $(f,g)$ is not primarily a
matrix pair; it is a point in the nilpotent cone distinguished from all other points
by the accessibility structure of its balanced and unbalanced sectors. An immersed
surface cobordism is not primarily a geometric shape; it is a movie equivalence class
under the rewrite grammar of extended Carter--Saito moves.

In each case, the \emph{object} emerges from the \emph{transformation algebra} as a
stable equivalence class, and the \emph{invariants} --- SAF class, nilpotency degree,
Khovanov homology --- are exactly what survives the identification under admissible
moves. This is what we mean by the claim that constraint-first systems do not derive
coherence from objects; rather, objects emerge as stable equivalence classes of
admissible transformations.

\subsection{Obstruction sensitivity as a recurring invariant principle}

Each theory also shares what we term \emph{obstruction sensitivity}: the invariants of
the system encode not positive information about configurations but negative information
about which transformations cannot be performed. The SAF class of a weighted foam
records what cannot be undone by any sequence of elementary cobordisms. The nilpotency
degree records how many round-trips are required to exhaust all accessible structure.
The Khovanov homology records which linking and knotting data cannot be eliminated by
any concordance.

This is characteristic of all genuine topological and homological invariants: they are
not descriptive of positive structure but \emph{obstructive} of positive transformation.
A class in $H_1(\AutIET, \ZZ)$ obstructs the lifting of a transformation to a product
of commutators. A nonzero element in $\Kh(L)$ obstructs the existence of a null-cobordism
of $L$ in the ball. A nonzero nilpotency degree obstructs the existence of a nonzero
vector surviving all iterated round-trips.

Obstruction sensitivity is the mathematical signature of constraint-first geometry: the
invariants are defined by what the constraint system excludes rather than by what it
permits.

\subsection{Formal introduction of CLIO}

We now crystallize the operational concept that has been guiding the interpretations
throughout this essay into a formal definition. The Constraint-Leveraged Inference
Operator has appeared informally in each domain: as the compression map from foam
configurations to SAF classes, as the nilpotency filter selecting admissible transport
pairs, and as the chart projection encoding immersed surfaces in reduced dimension.
Its formal role is that of a constrained projection preserving admissibility invariants
under trajectory compression.

\begin{definition}[CLIO --- Constraint-Leveraged Inference Operator]
\label{def:CLIO}

Let $\calX$ be a state or configuration space, $\calA$ a system of local admissibility
constraints (compatibility laws, cobordism moves, nilpotency conditions, movie-move
equivalences), and $\calM$ a compressed invariant space. A \emph{Constraint-Leveraged
Inference Operator} is a map
\[
\CLIO_{\calA} : \calX \to \calM
\]
satisfying the following conditions:

\begin{enumerate}

\item \emph{Admissibility preservation.}
$\CLIO_{\calA}$ is constant on $\calA$-equivalence classes: if
$x \sim_{\calA} x'$ then
\[
\CLIO_{\calA}(x)=\CLIO_{\calA}(x').
\]

\item \emph{Obstruction faithfulness.}
The fibers of $\CLIO_{\calA}$ are exactly the
$\calA$-equivalence classes:
\[
\CLIO_{\calA}(x)=\CLIO_{\calA}(x')
\quad \text{if and only if} \quad
x \sim_{\calA} x'.
\]

\item \emph{Singular extension.}
$\CLIO_{\calA}$ extends to a class of controlled
singular inputs
\[
\hat{\calX} \supset \calX
\]
in which singularities are treated as structured morphism types carrying definite
invariant data rather than as excluded configurations.

\end{enumerate}

\end{definition}

This definition is deliberately abstract. Its instances are the concrete maps studied
in the three source papers: $\psi : \Cob^1_{>0} \to \RR \wedgeQ \RR$ in foam theory
(Theorem~\ref{thm:IK-main}), the balanced-vector bijection of
Theorem~\ref{thm:CIKLR-main} in nilpotent transport theory, and the extended 2-functor
$\tilde\kappa$ in immersed cobordism theory.

Each instance satisfies the three conditions of Definition~\ref{def:CLIO}. In foam
theory: $\psi$ is constant on cobordism classes by construction, is an isomorphism
(hence obstruction-faithful), and the planar unoriented extension to $Z^2(\RR_{>0})$
provides the singular extension via the antisymmetric bracket. In nilpotent transport
theory: the bijection of Theorem~\ref{thm:CIKLR-main} identifies the nilpotent cone
as the admissibility-preserved image, the balanced/unbalanced dichotomy is
obstruction-faithful, and the Boolean semiring extension provides singular extension
via DAG structure. In immersed cobordism theory: $\tilde\kappa$ is constant on isotopy
classes by the movie-move theorem, is well-defined on homotopy classes of chain maps,
and the singular extension condition is exactly the addition of node morphism types
$A$ and $B$.

\subsection{The RSVP framework as admissibility meta-language}

We close with a remark on the role of RSVP in this essay. The Relativistic Scalar-Vector
Plenum framework has appeared throughout not as a primary mathematical object seeking
validation from topology but as a \emph{meta-interpretive language} for recognizing
admissibility geometry across independent mathematical systems. RSVP provides the
vocabulary for naming what the three source papers independently formalize: process-
primary ontology, obstruction-sensitive invariants, and admissibility-governed coherence.

The RSVP framework's central claim --- that physical and semantic coherence arises from
admissible trajectory equivalence rather than from intrinsic object properties --- is
precisely the mathematical content of Theorem~\ref{thm:IK-main}, Theorem~\ref{thm:CIKLR-main},
and the main theorem of \cite{CCKK25}, each formulated in their own domain. The present
essay argues that this convergence is not coincidental. The three papers are
independently discovering the same constraint-first geometry because that geometry is a
fundamental feature of any mathematical system in which coherence is preserved by
transformation rather than configuration.

\begin{strucprin}[Constraint-First Process Geometry]
\label{sp:main}
Coherent structure in a constraint-first system is not a property of configurations
but of transformation equivalence classes: objects emerge as stable equivalence classes
of admissible local operations, invariants encode obstruction to simplification under
admissible moves, and the admissibility geometry --- the system of local compatibility
laws --- is the primary datum from which all other structure is derived. The cobordism
group of weighted foams, the nilpotent cone of bidirectional transport pairs, and the
extended Khovanov 2-functor on immersed cobordisms are three independent mathematical
instantiations of this single architectural principle.
\end{strucprin}

% ============================================================
\section{Conclusion}
% ============================================================

We have traced a single architectural pattern through three domains of contemporary
mathematics: the foam cobordism theory of Im and Khovanov, the nilpotent pair theory of
Chen, Im, Khovanov, Lillja, and Rugo, and the immersed surface cobordism theory of
Carter, Cooper, Khovanov, and Krushkal. In each domain, the fundamental mathematical
entities are not configurations but equivalence classes of admissible local
transformations. Invariants are defined by obstruction to simplification. Singularities
carry structured algebraic information rather than representing failures of theory.
Global coherence is guaranteed not by metric embedding but by the consistent closure of
local compatibility laws.

The Constraint-Leveraged Inference Operator, introduced formally in
Definition~\ref{def:CLIO}, abstracts this common architecture: it is any map from a
configuration space to a compressed invariant space that is constant on admissibility
equivalence classes, faithful to the obstruction structure of those classes, and
extensible to controlled singular inputs carrying definite invariant data.

The RSVP framework provides the interpretive language for recognizing this architecture
across domains. Its central claim --- that physical and semantic coherence is an
admissibility-geometric phenomenon rather than an intrinsic object property --- is
independently confirmed by each of the three source papers. The convergence is not
imposed by external metaphysical framing but reflects a genuine structural feature of
constraint-governed mathematical systems.

The deepest implication is philosophical as much as mathematical. A constraint-first
ontology inverts the usual direction of explanation. It does not ask: ``Given objects
with intrinsic properties, which transformations are admissible?'' It asks instead:
``Given admissibility laws, which objects are stably distinguishable?'' The answer ---
the orbits, invariants, and equivalence classes we have examined --- constitutes not
merely a different calculation but a fundamentally different conception of what
mathematical structure is.

Constraint-first systems do not derive coherence from objects. Objects emerge as stable
equivalence classes of admissible transformations.

% ============================================================
% Bibliography placeholder
% ============================================================
\newpage
\begin{thebibliography}{99}

% ---- Primary sources ----

\bibitem{IK24foam}
Mee Seong Im and Mikhail Khovanov,
\textit{Foam cobordism and the Sah--Arnoux--Fathi invariant},
preprint, April 2024.

\bibitem{CIKLR25}
Weixi Chen, Mee Seong Im, Mikhail Khovanov, Catherine Lillja, and Nicolas Rugo,
\textit{Pairs of eventually constant maps and nilpotent pairs},
preprint, December 2025.

\bibitem{CCKK25}
Scott Carter, Benjamin Cooper, Mikhail Khovanov, and Vyacheslav Krushkal,
\textit{An extension of Khovanov homology to immersed surface cobordisms},
arXiv:2510.14760, October 2025.

% ---- K-theory and assemblers ----

\bibitem{Zak17a}
Inna Zakharevich,
\textit{The K-theory of assemblers},
\textit{Advances in Mathematics} \textbf{304} (2017), 1176--1218.

\bibitem{Zak17b}
Inna Zakharevich,
\textit{On $K_1$ of an assembler},
\textit{Journal of Pure and Applied Algebra} \textbf{221} (2017), no.~7, 1867--1898.

\bibitem{GIKK2024}
David Gepner, Mee Seong Im, Mikhail Khovanov, and Nitu Kitchloo,
\textit{Foam cobordisms and algebraic K-theory},
preprint, 2024.

\bibitem{Quillen1973}
Daniel Quillen,
\textit{Higher algebraic K-theory I},
Lecture Notes in Mathematics, vol.~341,
Springer, 1973.

% ---- Interval exchange transformations and SAF invariant ----

\bibitem{Vee84}
William A.~Veech,
\textit{The metric theory of interval exchange transformations III:
The Sah--Arnoux--Fathi invariant},
\textit{American Journal of Mathematics} \textbf{106} (1984), no.~6, 1389--1422.

\bibitem{Vorobets2011}
Yaroslav Vorobets,
\textit{Notes on the Sah--Arnoux--Fathi invariant},
preprint, 2011.

\bibitem{DS2016}
Vincent Delecroix and Pascal Hubert,
\textit{Sah--Arnoux--Fathi invariant and interval exchange transformations},
\textit{Ergodic Theory and Dynamical Systems} (2016).

\bibitem{Arnoux1981}
Pierre Arnoux,
\textit{{\'E}changes d'intervalles et flots sur les surfaces},
in \textit{Ergodic Theory and Dynamical Systems}, 1981.

\bibitem{Zorich2006a}
Anton Zorich,
\textit{Flat surfaces},
in \textit{Frontiers in Number Theory, Physics, and Geometry~I},
Springer, 2006.

\bibitem{Zorich2006b}
Anton Zorich,
\textit{Deviation for interval exchange transformations},
\textit{Ergodic Theory and Dynamical Systems} (2006).

% ---- Khovanov homology and foams ----

\bibitem{Kho00}
Mikhail Khovanov,
\textit{A categorification of the Jones polynomial},
\textit{Duke Mathematical Journal} \textbf{101} (2000), no.~3, 359--426.

\bibitem{Kho03}
Mikhail Khovanov,
\textit{Patterns in knot cohomology~I},
\textit{Experimental Mathematics} \textbf{12} (2003), no.~3, 365--374.

\bibitem{Kho06a}
Mikhail Khovanov,
\textit{An invariant of tangle cobordisms},
\textit{Transactions of the American Mathematical Society} \textbf{358} (2006), no.~1, 315--327.

\bibitem{Kho06b}
Mikhail Khovanov,
\textit{Link homology and Frobenius extensions},
\textit{Fundamenta Mathematicae} \textbf{190} (2006), 179--190.

\bibitem{Kho04}
Mikhail Khovanov,
\textit{sl(3) link homology},
\textit{Algebraic \& Geometric Topology} \textbf{4} (2004), 1045--1081.

\bibitem{BN05}
Dror Bar-Natan,
\textit{Khovanov's homology for tangles and cobordisms},
\textit{Geometry and Topology} \textbf{9} (2005), 1443--1499.

\bibitem{BN07}
Dror Bar-Natan,
\textit{Fast Khovanov homology computations},
\textit{Journal of Knot Theory and Its Ramifications} \textbf{16} (2007), no.~3, 243--255.

\bibitem{MV2007}
Marco Mackaay and Pedro Vaz,
\textit{The universal sl(3)-link homology},
\textit{Algebraic \& Geometric Topology} \textbf{7} (2007), 1135--1169.

\bibitem{RW2016}
Daniel Rose and Paul Wedrich,
\textit{Deformations of colored sl(N) link homologies via foams},
\textit{Geometry \& Topology} (2016).

\bibitem{RW2020a}
Daniel Rose and Paul Wedrich,
\textit{Categorified quantum link invariants and foams},
\textit{Advances in Mathematics} (2020).

\bibitem{KM2019}
Peter Kronheimer and Tomasz Mrowka,
\textit{Instanton Floer homology and orbifolds},
\textit{Geometry \& Topology} (2019).

\bibitem{QW22}
Hoel Queffelec and Paul Wedrich,
\textit{Movie moves for framed foams from multijet transversality},
arXiv:2205.14947, 2022.

% ---- Surface cobordisms in 4-space ----

\bibitem{CS98}
J.~Scott Carter and Masahico Saito,
\textit{Knotted surfaces and their diagrams},
Mathematical Surveys and Monographs, vol.~55,
American Mathematical Society, Providence, RI, 1998.

\bibitem{CRS97}
J.~Scott Carter, J{\"o}rg Rieger, and Masahico Saito,
\textit{A combinatorial description of knotted surfaces and their isotopies},
\textit{Advances in Mathematics} \textbf{127} (1997), 1--51.

\bibitem{Jac04}
Magnus Jacobsson,
\textit{An invariant of link cobordisms from Khovanov homology},
\textit{Algebraic \& Geometric Topology} \textbf{4} (2004), 1211--1251.

\bibitem{CMW09}
David Clark, Scott Morrison, and Kevin Walker,
\textit{Fixing the functoriality of Khovanov homology},
\textit{Geometry \& Topology} \textbf{13} (2009), 1499--1582.

\bibitem{San21}
Tetsuya Sano,
\textit{Functoriality of Khovanov homology revisited},
\textit{Journal of Knot Theory and Its Ramifications} \textbf{30} (2021), no.~14.

\bibitem{Bla10}
Christian Blanchet,
\textit{An oriented model for Khovanov homology},
\textit{Journal of Knot Theory and Its Ramifications} \textbf{19} (2010), no.~2, 291--312.

\bibitem{Ras10}
Jacob Rasmussen,
\textit{Khovanov homology and the slice genus},
\textit{Inventiones Mathematicae} \textbf{182} (2010), 419--447.

\bibitem{MWW22}
Scott Morrison, Kevin Walker, and Paul Wedrich,
\textit{Invariants of 4-manifolds from Khovanov--Rozansky link homology},
\textit{Geometry \& Topology} \textbf{26} (2022), no.~8, 3367--3420.

\bibitem{RW24}
Qiuyu Ren and Michael Willis,
\textit{Khovanov homology and exotic 4-manifolds},
arXiv:2402.10452, 2024.

\bibitem{Don83}
Simon Donaldson,
\textit{An application of gauge theory to four-dimensional topology},
\textit{Journal of Differential Geometry} \textbf{18} (1983), 279--315.

\bibitem{Fre82}
Michael Freedman,
\textit{The topology of four-dimensional manifolds},
\textit{Journal of Differential Geometry} \textbf{17} (1982), 357--453.

\bibitem{Kro97}
Peter Kronheimer,
\textit{An obstruction to removing intersection points in immersed surfaces},
\textit{Topology} \textbf{36} (1997), no.~4, 931--962.

\bibitem{ISST25}
Hayato Imori, Taketo Sano, Kouki Sato, and Masaki Taniguchi,
\textit{Cobordism maps in Khovanov homology and singular instanton homology~II},
arXiv:2510.09399, 2025.

\bibitem{ItoYoshida21}
Noboru Ito and Jun Yoshida,
\textit{A cobordism realizing crossing change on sl2 tangle homology and a
categorified Vassiliev skein relation},
\textit{Topology and Its Applications} \textbf{296} (2021).

\bibitem{ItoYoshida23}
Noboru Ito and Jun Yoshida,
\textit{On the four-term relation on Khovanov homology},
arXiv:2202.08527, 2023.

\bibitem{Car15}
J.~Scott Carter,
\textit{Reidemeister/Roseman-type moves to embedded foams in 4-dimensional space},
in \textit{New Ideas in Low Dimensional Topology},
World Scientific, 2015.

\bibitem{CK12}
J.~Scott Carter and Seiichi Kamada,
\textit{Diagrammatic algebra},
Mathematical Surveys and Monographs, vol.~264,
American Mathematical Society, 2021.

\bibitem{CH15}
Benjamin Cooper and Matt Hogancamp,
\textit{An exceptional collection for Khovanov homology},
\textit{Algebraic \& Geometric Topology} \textbf{15} (2015), 2659--2707.

\bibitem{Cap08}
Carmen Caprau,
\textit{sl(2) tangle homology with a parameter and singular cobordisms},
\textit{Algebraic \& Geometric Topology} \textbf{8} (2008), 729--756.

\bibitem{Wen11}
Luoying Weng,
\textit{Isotopy invariants of immersed surfaces in a 4-manifold},
Ph.D.\ thesis, State University of New York at Stony Brook, 2011.

\bibitem{Kau89}
Louis H.~Kauffman,
\textit{Invariants of graphs in three-space},
\textit{Transactions of the American Mathematical Society} \textbf{311} (1989), no.~2, 697--710.

\bibitem{Kau13}
Louis H.~Kauffman,
\textit{Knots and physics},
4th edition, World Scientific, 2013.

\bibitem{BK01}
Bojko Bakalov and Alexander Kirillov Jr.,
\textit{Lectures on tensor categories and modular functors},
University Lecture Series, vol.~21,
American Mathematical Society, 2001.

% ---- Train tracks, measured laminations, branched surfaces ----

\bibitem{PH92}
Robert C.~Penner and John L.~Harer,
\textit{Combinatorics of train tracks},
Annals of Mathematics Studies, vol.~125,
Princeton University Press, Princeton, NJ, 1992.

\bibitem{Oertel1984}
Ulrich Oertel,
\textit{Incompressible branched surfaces},
\textit{Inventiones Mathematicae} \textbf{76} (1984), 385--410.

\bibitem{Oertel1988}
Ulrich Oertel,
\textit{Measured laminations in 3-manifolds},
\textit{Transactions of the American Mathematical Society} \textbf{305} (1988), no.~2, 531--573.

\bibitem{Thurston1988}
William P.~Thurston,
\textit{On the geometry and dynamics of diffeomorphisms of surfaces},
\textit{Bulletin of the American Mathematical Society} \textbf{19} (1988), no.~2, 417--431.

% ---- Nilpotent operators, quiver representations, combinatorics ----

\bibitem{Lei21}
Tom Leinster,
\textit{The probability that an operator is nilpotent},
\textit{American Mathematical Monthly} \textbf{128} (2021), no.~4, 371--375.

\bibitem{FH58}
N.~J.~Fine and I.~N.~Herstein,
\textit{The probability that a matrix be nilpotent},
\textit{Illinois Journal of Mathematics} \textbf{2} (1958), 499--504.

\bibitem{CG10}
Neil Chriss and Victor Ginzburg,
\textit{Representation theory and complex geometry},
Birkh{\"a}user Boston, 2010.

\bibitem{GR25}
Lydia G{\"o}smann and Markus Reineke,
\textit{Motives of nullcones of quiver representations},
arXiv:2502.02353, 2025.

\bibitem{QuiverRep}
William Crawley-Boevey,
\textit{Lectures on representations of quivers},
available online.

\bibitem{Gabriel}
Pierre Gabriel,
\textit{Unzerlegbare Darstellungen~I},
\textit{Manuscripta Mathematica} \textbf{6} (1972), 71--103.

\bibitem{AS90}
Moh'd Z.~Abu-Sbeih,
\textit{On the number of spanning trees of $K_n$ and $K_{m,n}$},
\textit{Discrete Mathematics} \textbf{84} (1990), 205--207.

\bibitem{FS58}
Miroslav Fiedler and Ji\v{r}{\'{\i}} Sedl{\'a}\v{c}ek,
\textit{On $w$-bases of directed graphs},
\v{C}asopis pro p\v{e}stov{\'a}n{\'{\i}} matematiky \textbf{83} (1958), 214--225.

\bibitem{Joy81}
Andr{\'e} Joyal,
\textit{Une th{\'e}orie combinatoire des s{\'e}ries formelles},
\textit{Advances in Mathematics} \textbf{42} (1981), 1--82.

\bibitem{But72}
Kim Ki-Hang Butler,
\textit{The number of idempotents in $(0,1)$-matrix semigroups},
\textit{Linear Algebra and its Applications} \textbf{5} (1972), no.~3, 233--246.

\bibitem{DK25}
Nayana Shibu Deepthi and Chanchal Kumar,
\textit{Combinatorial identities using the matrix tree theorem},
arXiv:2504.21319, 2025.

\bibitem{IK22}
Mee Seong Im and Mikhail Khovanov,
\textit{Topological theories and automata},
arXiv:2202.13398, 2022.

\bibitem{Jac89}
Nathan Jacobson,
\textit{Basic algebra~II},
second edition,
W.~H.~Freeman and Company, New York, 1989.

% ---- Category theory and topology ----

\bibitem{MacLane}
Saunders Mac Lane,
\textit{Categories for the working mathematician},
second edition,
Graduate Texts in Mathematics 5,
Springer, 1998.

\bibitem{Sheaves}
Masaki Kashiwara and Pierre Schapira,
\textit{Categories and sheaves},
Springer, 2006.

\bibitem{LurieHTT}
Jacob Lurie,
\textit{Higher topos theory},
Princeton University Press, 2009.

\bibitem{LurieHA}
Jacob Lurie,
\textit{Higher algebra},
preprint.

\bibitem{Grothendieck1957}
Alexander Grothendieck,
\textit{Sur quelques points d'alg\`ebre homologique},
\textit{T\^{o}hoku Mathematical Journal} (1957).

\bibitem{BaezStay}
John Baez and Mike Stay,
\textit{Physics, topology, logic and computation: a Rosetta Stone},
in \textit{New Structures for Physics},
Springer, 2010.

\bibitem{Mun2000}
James R.~Munkres,
\textit{Topology},
second edition,
Prentice Hall, 2000.

\bibitem{Milnor}
John Milnor,
\textit{Morse theory},
Princeton University Press, 1963.

\bibitem{Gromov}
Mikhail Gromov,
\textit{Partial differential relations},
Springer, 1986.

% ---- Thermodynamics, entropy, and physics ----

\bibitem{Jacobson}
Ted Jacobson,
\textit{Thermodynamics of spacetime: the Einstein equation of state},
\textit{Physical Review Letters} \textbf{75} (1995), no.~7, 1260--1263.

\bibitem{Verlinde}
Erik Verlinde,
\textit{On the origin of gravity and the laws of Newton},
\textit{Journal of High Energy Physics} \textbf{2011} (2011), no.~4, 29.

\bibitem{EntropyGeometry}
Carlo Rovelli,
\textit{Statistical mechanics of gravity and the thermodynamical origin of time},
\textit{Classical and Quantum Gravity} \textbf{10} (1993), 1549--1566.

\bibitem{Barandes}
Jacob Barandes,
\textit{A suggestion for quantum theory in terms of unistochastic matrices},
arXiv preprint, 2023.

% ---- Cognitive and semantic frameworks ----

\bibitem{cheng2025cognitiveloopinsituoptimization}
Newman Cheng, Gordon Broadbent, and William Chappell,
\textit{Cognitive loop via in-situ optimization: self-adaptive reasoning for science},
arXiv:2508.02789, 2025.

\bibitem{Friston}
Karl Friston,
\textit{The free-energy principle: a unified brain theory?}
\textit{Nature Reviews Neuroscience} \textbf{11} (2010), 127--138.

\bibitem{BM76}
John A.~Bondy and Uppaluri S.~R.~Murty,
\textit{Graph theory with applications},
American Elsevier Publishing Company, 1976.

\bibitem{CCPS11}
William J.~Cook, William H.~Cunningham, William R.~Pulleyblank, and Alexander Schrijver,
\textit{Combinatorial optimization},
Wiley Series in Discrete Mathematics and Optimization, vol.~33,
John Wiley \& Sons, 2011.

\end{thebibliography}

\end{document}
